\documentclass[a4paper,10pt]{article}
\usepackage[francais]{babel}
\usepackage{url}

\selectlanguage{francais}
\input{letter-icube-uds.sty}

\selectlanguage{francais}
\begin{document}


\name{St\'ephane \textsc{Genaud}}
\fonction{Ma\^itre de conf\'erences (HDR)}
\telephoneillkirch{+33 (0)3688 54542}
\email{genaud@unistra.fr}

\begin{letter}[Recommendation de Vincent Kherbache pour une thèse de doctorat]{à}{Commission de sélection,\\INRIA}

Madame, Monsieur

Je connais Vincent depuis son TER (Travail d'Etudes et Recherche) que j'ai
encadré (janvier-juin 2012) lors de sa première année de Master en informatique,
spécialité Réseaux et Systèmes embarqués. 
Le travail confié à Vincent dans ce TER était la mise en place d'un environnement
logiciel permettant de contrôler l'allocation de jobs à des machines virtuelles (VM)
gérées par un gestionnaire de VM. En l'occurence, nous souhaitions pouvoir contrôler
avec le batch-scheduler Slurm, l'allocation de jobs sur des VMs gérées avec Eucalyptus. 
Vincent a rempli sa mission avec succès en faisant preuve de beaucoup d'autonomie. 
J'ai particulièrement apprécié chez Vincent la qualité de son écoute. Il a, à mon
avis, beaucoup progressé en tenant scrupuleusement compte des conseils que nous lui
donnions. La présentation finale de son travail de TER a été très bien construite et 
documentée.


Vincent a souhaité poursuivre sur le sujet avec un stage d'été (facultatif) et nous
l'avons accueilli avec plaisir pour trois mois. Il a travaillé sur Schlouder, un ordonnanceur
côté utilisateur, qui s'appuie sur l'environment logiciel réalisé pendant son TER.
Cette fois, le travail concerne beaucoup plus les aspects algorithmiques. Il
nécessite cependant de toujours faire de bon choix systèmes pour 
que l'ordonnanceur puisse communiquer avec les composants Slurm et Eucalyptus.


Durant ce stage, Vincent a travaillé en équipe avec un autre étudiant de niveau
comparable, et avec un doctorant. Vincent a montré lors de ce travail qu'il savait 
travailler en équipe.  Ceci n'est pas surprenant dans la mesure où Vincent est très 
ouvert et constamment à l'écoute des autres.
Au final, le résultat du travail a été très fructueux. 


Enfin, Vincent a souhaité effectuer son stage de recherche Master dans notre équipe.
Initialement intéréssé par la sécurité dans les clouds, il s'est résolu à 
travailler sur un sujet dans notre domaine de compétence. Il travaille actuellement
sur les stratégies de provisioning et d'ordonnancement de ressources dans un
contexte multi-clouds. Il a commencé par la modularisation du composant gérant
les VMs, en faisant que le système fonctionne aussi bien avec OpenStack qu'Eucalyptus.
Il aborde maintenant des stratégies multi-objectifs impliquant potentiellement plusieurs
clouds.


Du point de vue des résultats scolaires, il faut également remarquer que Vincent 
a obtenu de très bonnes notes aux épreuves théoriques de son Master, et se classe,
au premier semestre (avant le stage de recherche), 2e sur 19 étudiants.

En conclusion, je recommande vivement Vincent pour une thèse de doctorat, à la 
fois pour ses qualités humaines et son sérieux dans le travail. 


\end{letter}
\begin{flushright}
\includegraphics[width=.20\textwidth]{signgenaud.jpg}
\end{flushright}
%\newpage
%\bibliographystyle{plain}
%\bibliography{julien_blaise_ED}


\end{document}
